\documentclass[12pt]{article}

\usepackage[utf8]{inputenc}
\usepackage{amsmath, amssymb, amsthm}
\usepackage{geometry}
\geometry{a4paper, margin=2.5cm}

\title{Teorema lui Schwarz}
\author{}
\date{}

\theoremstyle{definition}
\newtheorem{definition}{Definiție}

\theoremstyle{plain}
\newtheorem{theorem}{Teoremă}
\newtheorem{proposition}{Propoziție}

\begin{document}

\maketitle

\begin{definition}
Fie $u : \Omega \subset \mathbb{R}^2 \to \mathbb{R}$ o funcție.  
Presupunem că $u$ admite derivate parțiale de ordinul doi în punctul $(x_0, y_0)$, adică


\[
u_{xy}(x_0,y_0), \qquad u_{yx}(x_0,y_0)
\]


există. Acestea se numesc \emph{derivate parțiale încrucișate} ale lui $u$ în punctul $(x_0,y_0)$.
\end{definition}

\begin{theorem}
(\textbf{Teorema lui Schwarz})  
Fie $u : \Omega \subset \mathbb{R}^2 \to \mathbb{R}$ o funcție de clasă $C^2$ pe $\Omega$, adică toate derivatele parțiale de ordinul doi ale lui $u$ există și sunt continue pe $\Omega$.  
Atunci, pentru orice punct $(x,y) \in \Omega$, are loc egalitatea


\[
u_{xy}(x,y) = u_{yx}(x,y).
\]


Cu alte cuvinte, derivatele parțiale încrucișate sunt egale.
\end{theorem}

\begin{proposition}
(\textbf{Consecință})  
Dacă $u$ este de clasă $C^2$ pe $\Omega$, atunci matricea Hessiană


\[
H_u(x,y) =
\begin{pmatrix}
u_{xx}(x,y) & u_{xy}(x,y) \\
u_{yx}(x,y) & u_{yy}(x,y)
\end{pmatrix}
\]


este simetrică în orice punct $(x,y) \in \Omega$.
\end{proposition}

\end{document}
